\documentclass{article}
\usepackage{amsmath, amssymb, amsthm}
\usepackage[margin=1in]{geometry}
\usepackage{hyperref}

\newtheorem{theorem}{Theorem}
\newtheorem{lemma}[theorem]{Lemma}
\newtheorem{proposition}[theorem]{Proposition}
\newtheorem{corollary}[theorem]{Corollary}
\newtheorem{conjecture}[theorem]{Conjecture}
\theoremstyle{definition}
\newtheorem{definition}[theorem]{Definition}
\newtheorem{remark}[theorem]{Remark}

\title{$\sigma_2$-G1 at $n=6$: \\ KL-cone Decomposition via the Doubled Word}
\author{Clio}
\date{2026-04-29}

\begin{document}

\maketitle

\begin{abstract}
We refine the trace-square reduction of $\sigma_2$-G1 by decomposing both
$\sigma_1^2$ and $\operatorname{tr}(\Pi_q^2)$ in the unique linear basis of
$\mathbb{Z}[q]$ given by $\{(1+q)^l q^{(M-l)/2}\}_{l\,\text{even},\,0\leq l \leq M}$,
where $M = 2N$ and $N = \binom{n}{2}$. Each decomposition has non-negative
coefficients by $\sigma_1$-G1 (applied to the staircase word and to its
doubling, respectively). The strong-G1 conjecture
$N_l(\lambda) \geq m_l(\lambda)$ would then place $\sigma_2$ in the positive
cone of basis vectors, immediately yielding palindromic positivity.
We verify this strong-G1 cone bound \emph{computationally} for
$\lambda = (5,1)$ and $\lambda = (4,2)$ at $n=6$, but \emph{refute} it for
$\lambda = (3,2,1)$ at $n=6$ (negative coefficient $-3$ at $l=4$). Thus
KL-cone decomposition closes $\sigma_2$-G1 at $n=6$ for two of three rank-$\geq 2$
shapes; $(3,2,1)$ requires the orthogonal atom-decomposition of
\texttt{2026-04-29-sigma2-G1-321-atom-decomp.tex}. Together the two methods
close $\sigma_2$-G1 at $n=6$ in full, but not via a single uniform mechanism.
\end{abstract}

\section{Setup}

Let $H_q = H_q(S_n)$ be the Iwahori--Hecke algebra with quadratic relation $T_i^2 = (q-1)T_i + q$.
Let $\Pi_q^{(n)} = \prod_{j=1}^N (1+T_{i_j})$ be the Bott--Samelson operator along the
staircase reduced word for $w_0$, with $N = \binom{n}{2}$, and let $V_\lambda$ denote
Hoefsmit's seminormal Specht module for $\lambda \vdash n$. Write
$\sigma_k(V_\lambda) := e_k(\text{eigenvalues of } \Pi_q^{(n)}|_{V_\lambda})$.

\begin{definition}[$\sigma_k$-G1]
A polynomial $f(q) \in \mathbb{Z}[q]$ is \emph{G1} if
$f(q) = q^a (1+q)^b \cdot P(q)$ with $P \in \mathbb{Z}_{\geq 0}[q]$ palindromic.
We say $\sigma_k$-G1 holds for $V_\lambda$ if $\sigma_k(V_\lambda)$ is G1.
\end{definition}

\begin{theorem}[$\sigma_1$-G1, Apr 24]\label{thm:sigma1G1}
For every $\lambda \vdash n$,
$\sigma_1(V_\lambda) = \operatorname{tr}(\Pi_q^{(n)}|V_\lambda) = \sum_{k=0}^N n_k(\lambda)\,(1+q)^k\, q^{(N-k)/2}$
with $n_k(\lambda) \in \mathbb{Z}_{\geq 0}$, where the sum is over $k$ with $N-k$ even.
\end{theorem}

The proof of Theorem~\ref{thm:sigma1G1} uses the Kazhdan--Lusztig decomposition
$T_i + 1 = (1+q)D_i + vM_i$ (with $v = q^{1/2}$, $D_i$ a diagonal indicator,
and $M_i$ a non-negative integer matrix). Expanding the BS product and applying
the parity-vanishing lemma yields the stated form.

\section{Doubled-word analogue and basis decomposition}

\begin{lemma}[Doubled-word $\sigma_1$-G1]\label{lem:doubled}
For every $\lambda \vdash n$,
\[
  \operatorname{tr}((\Pi_q^{(n)})^2 | V_\lambda) \;=\; \sum_{l=0}^{2N} m_l(\lambda)\,(1+q)^l\, q^{(2N-l)/2},
\]
with $m_l(\lambda) \in \mathbb{Z}_{\geq 0}$ and $m_l = 0$ unless $l$ is even.
\end{lemma}

\begin{proof}
The proof of Theorem~\ref{thm:sigma1G1} uses only that $\Pi_q^{(n)}$ is a product of
factors $(1+T_{i_j})$, applying the KL decomposition $T_i+1 = (1+q)D_i + vM_i$
factorwise and then taking the trace with the parity vanishing lemma.
The same proof applies verbatim to any product of factors $(1+T_{i_j})$ — in
particular to the doubled BS-word $\Pi_q^{(n)} \cdot \Pi_q^{(n)}$, viewed as
$\prod_{j=1}^{2N}(1+T_{i_{(j \bmod N)}})$ along the doubled staircase index sequence.
The only change is $N \mapsto 2N$.
\end{proof}

\begin{proposition}[Basis $B_M$]\label{prop:basis}
The set $B_M := \{(1+q)^l q^{(M-l)/2} : l \in \{0, 2, 4, \dots, M\}\}$ is a
$\mathbb{Q}$-linear basis of the subspace of $\mathbb{Q}[q]$ spanned by
$\{q^a : 0 \leq a \leq M\}$ that consists of polynomials of degree $\leq M$
which are palindromic about $q^{M/2}$ (i.e., $f(q) = q^M f(1/q)$). Both
$\sigma_1^2$ and $\operatorname{tr}(\Pi_q^2)$ lie in this subspace with $M = 2N$.
\end{proposition}

\begin{proof}
The basis vectors are linearly independent: each has distinct minimum degree
$(M-l)/2$. There are $M/2 + 1$ basis vectors and $M/2 + 1$ degrees of freedom
in the palindromic-about-$q^{M/2}$ subspace of degree $\leq M$ polynomials.
$\sigma_1$ is palindromic about $q^{N/2}$ (eigenvalues of $\Pi_q$ are palindromic
in $q$ by parity vanishing), so $\sigma_1^2$ is palindromic about $q^N$.
$\operatorname{tr}(\Pi_q^2)$ is similarly palindromic about $q^N$ by Lemma~\ref{lem:doubled}.
\end{proof}

\section{KL-cone decomposition of $\sigma_2$}

\begin{theorem}[$\sigma_2$ in $B_M$]\label{thm:sigma2-decomp}
With $M = 2N$,
\[
  \sigma_1^2(V_\lambda) = \sum_{l \text{ even}} N_l(\lambda)\, (1+q)^l\, q^{(M-l)/2}, \qquad
  N_l(\lambda) = \sum_{k+k'=l} n_k(\lambda)\, n_{k'}(\lambda) \;\geq\; 0,
\]
\[
  \operatorname{tr}(\Pi_q^2 | V_\lambda) = \sum_{l \text{ even}} m_l(\lambda)\, (1+q)^l\, q^{(M-l)/2}, \qquad m_l(\lambda) \geq 0,
\]
and consequently, by Newton's identity $2\sigma_2 = \sigma_1^2 - \operatorname{tr}(\Pi_q^2)$,
\[
  \sigma_2(V_\lambda) = \sum_{l \text{ even}} \frac{N_l(\lambda) - m_l(\lambda)}{2}\,(1+q)^l\, q^{(M-l)/2}.
\]
\end{theorem}

\begin{proof}
The first identity squares $\sigma_1 = \sum_k n_k (1+q)^k q^{(N-k)/2}$ and
collects monomials by total exponent of $(1+q)$. The second is Lemma~\ref{lem:doubled}.
The third is direct: by Proposition~\ref{prop:basis} the basis $B_M$ is linearly
independent, so coefficients combine.
\end{proof}

\begin{conjecture}[Strong-G1 cone bound]\label{conj:strong}
For every $\lambda \vdash n$ and every $l$,
\[
  N_l(\lambda) \;\geq\; m_l(\lambda),\qquad N_l(\lambda) - m_l(\lambda) \in 2\mathbb{Z}_{\geq 0}.
\]
\end{conjecture}

If Conjecture~\ref{conj:strong} holds, then by Theorem~\ref{thm:sigma2-decomp},
$\sigma_2$ lies in the positive integer cone spanned by $B_M$. Since each basis
vector $(1+q)^l q^{(M-l)/2}$ is itself G1 (with palindromic factor $1$),
$\sigma_2$ is then G1.

\section{Computational verification}

We computed $N_l, m_l$ for all $\lambda \vdash n$ with $n \leq 6$ and rank $\geq 1$
using the script \texttt{test\_KL\_decomposition\_sigma2.py}, which evaluates
$\Pi_q$ and $\Pi_q^2$ in Hoefsmit's seminormal basis and solves the linear
system from $B_M$ to extract coefficients.

\begin{table}[h]
\centering
\caption{$N_l - m_l$ for rank-$\geq 2$ shapes at $n=6$ ($N=15$, $M=30$).}
\begin{tabular}{r|rrrrrrrrrrrrrrr}
$l$ & 2 & 4 & 6 & 8 & 10 & 12 & 14 & 16 & 18 & 20 & 22 & 24 & 26 \\\hline
$(5,1)$ & 0 & 0 & 0 & 0 & 0 & 0 & 0 & 0 & 12 & 26 & 12 & 2 & 0 \\
$(4,2)$ & 0 & 0 & 0 & 0 & 84 & 194 & 160 & 78 & 22 & 2 & 0 & 0 & 0 \\
$(3,2,1)$ & 2 & $\boldsymbol{-6}$ & 18 & 24 & 12 & 2 & 0 & 0 & 0 & 0 & 0 & 0 & 0
\end{tabular}
\end{table}

For $(5,1)$ and $(4,2)$, all $N_l - m_l \geq 0$; Conjecture~\ref{conj:strong}
holds for these two shapes. The resulting $\sigma_2$ decompositions are:
\begin{align*}
\sigma_2(V_{(5,1)}) &= 6\,(1+q)^{18}q^6 + 13\,(1+q)^{20}q^5 + 6\,(1+q)^{22}q^4 + (1+q)^{24}q^3, \\
\sigma_2(V_{(4,2)}) &= 42\,(1+q)^{10}q^{10} + 97\,(1+q)^{12}q^9 + 80\,(1+q)^{14}q^8 \\
&\quad + 39\,(1+q)^{16}q^7 + 11\,(1+q)^{18}q^6 + (1+q)^{20}q^5,
\end{align*}
manifestly G1.

For $(3,2,1)$, $N_4 - m_4 = -6$, so the cone bound fails. Direct verification
gives
\[
  \sigma_2(V_{(3,2,1)}) = q^9(1+q)^2 \cdot P_{(3,2,1)}(q),
\]
with $P_{(3,2,1)}(q) = q^{10}+16q^9+105q^8+369q^7+759q^6+961q^5+759q^4+369q^3+105q^2+16q+1$
palindromic and $\geq 0$, so $\sigma_2$-G1 \emph{still holds}, just not via a
positive-cone decomposition in $B_M$.

\section{Why $(3,2,1)$ resists this approach}

The basis $B_M$ at $M = 30$ includes vectors with low $l$ and high $q$-power
(e.g., $l=2$ has support $[14,16]$). For shapes with large $b_2 = b_{\lambda,2}$
(the $(1+q)$ exponent in $\sigma_2$), the unique decomposition is forced into
the high-$l$ portion of $B_M$, leaving zero coefficient at low $l$. This is
the case for $(5,1)$ ($b_2 = 18$) and $(4,2)$ ($b_2 = 10$).

For $(3,2,1)$, $b_2 = 2$ is small. The $\sigma_2$ polynomial sits with leading
$q$-degree $9$ and trailing $q$-degree $21$, but its palindromic factor of
degree $12$ is "wide" in degree compared with $(1+q)^2$. The unique decomposition
in $B_M$ then uses low-$l$ basis vectors with negative coefficients to
``carve out'' the wide palindromic shape from the narrower $(1+q)^l$ profiles.

\begin{remark}[Interpretation]
The strong-G1 cone bound $N_l \geq m_l$ is genuinely stronger than $\sigma_2$-G1.
The failure at $(3,2,1)$ suggests that strong-G1 captures shapes whose
$\sigma_2$ has ``simple'' palindromic structure (factor $(1+q)^b$ dominating),
while $(3,2,1)$ has the genuinely $11$-coefficient palindromic atom $A_{10}$
that does not factor through $(1+q)$.
This matches the atom-alphabet picture from
\texttt{2026-04-29-sigma2-G1-321-atom-decomp.tex}: $(3,2,1)$ uses the
\emph{orphan} atom $A_{10} = q^{10}+16q^9+\dots$ (OEIS A344437, derangement
RTL-minima), unique among rank-$\geq 2$ shapes at $n \leq 6$.
\end{remark}

\section{Status and gap}

\paragraph{What this proves at $n=6$ (modulo Conjecture~\ref{conj:strong}).}
For $\lambda \in \{(5,1),\,(4,2)\}$, $\sigma_2$-G1 follows from the
KL-cone decomposition once $N_l \geq m_l$ is established.
\emph{Computationally} this is closed; \emph{structurally} it requires
proving the per-$l$ inequality
\[
  \sum_{k+k'=l} n_k(\lambda)\, n_{k'}(\lambda) \;\geq\; m_l(\lambda)
\]
combinatorially. This sharpens the gap from ``$\sigma_2$-G1 holds'' to a
specific cone-positivity question on KL-coefficients.

\paragraph{Combined closure of $\sigma_2$-G1 at $n=6$.}
\begin{enumerate}
\item $(5,1)$: KL-cone decomposition (this paper, modulo the cone bound, computationally verified).
\item $(4,2)$: KL-cone decomposition (this paper, modulo the cone bound, computationally verified).
\item $(3,2,1)$: orthogonal atom-alphabet decomposition (\texttt{2026-04-29-sigma2-G1-321-atom-decomp.tex}).
\end{enumerate}

This is a genuine \emph{partial} result: the lemma needed for $(5,1), (4,2)$
is the per-shape combinatorial inequality $N_l \geq m_l$, and a structural
proof of that inequality remains open. The $(3,2,1)$ case is closed by a
different mechanism.

\paragraph{Open: structural proof of the cone bound.}
A natural candidate is per-pair sub-multiplicativity:
\[
  \operatorname{tr}(P_S | V_\lambda) \cdot \operatorname{tr}(P_{S'} | V_\lambda) \;\geq\; \operatorname{tr}(P_S P_{S'} | V_\lambda),
\]
where $P_S$ is the $\sigma_1$-G1 operator product over the subset $S$. For
positive semi-definite $P_S, P_{S'}$ this would follow from the standard
inequality $\operatorname{tr}(AB) \leq \operatorname{tr}(A)\operatorname{tr}(B)$
for PSD operators. The KL operators $D_i, M_i$ are non-negative integer
matrices in Hoefsmit's basis, but their products are not always PSD on
$V_\lambda$, and the per-pair inequality fails for $(3,2,1)$ at $l=4$.

\paragraph{Why $\sigma_2$-G1 closes despite this.}
The strong-G1 cone bound is strictly stronger than $\sigma_2$-G1. The right
structural invariant for $\sigma_2$-G1 across all rank-$\geq 2$ shapes is
plausibly the atom alphabet, not the $B_M$ basis. The $B_M$ basis appears to
suffice precisely when the shape's palindromic factor is captured by KF
($q^a(1+q)^b$) elements; $(3,2,1)$ first introduces an irreducible
non-cyclotomic palindromic atom ($A_{10}$, see project memory
\texttt{project\_oeis\_derangement\_rtl.md}) and lies outside this cone.

\section*{Acknowledgements and references}

This refines the trace-square framing of
\texttt{2026-04-28-sigma2-G1-n6-trace-square.tex} by working in a fixed
linear basis. The atom-decomposition closure of $(3,2,1)$ is established in
\texttt{2026-04-29-sigma2-G1-321-atom-decomp.tex}. The original $\sigma_1$-G1
proof underpinning Lemma~\ref{lem:doubled} is in
\texttt{2026-04-24-sigma1-G1.tex}. Computational verification:
\texttt{/home/clio/projects/scratch/2026-04-29-s5-equivariance/test\_KL\_decomposition\_sigma2.py}.

\end{document}
